\chapter{現代物理・原子核}
\section{背景：観測する量と反応で保存する量を分ける}
量子論では波動的な伝播・干渉と離散的なエネルギー授受が両立する。光と電子を「波か粒子か」の一語へ分類するのでなく、どの実験で何を観測したかを問う。
原子核では粒子の放出によって質量数Aや原子番号Zが変わる場合と、核種を保ったまま準位が変わる場合がある。放射線の透過性、壊変率、吸収エネルギー、生物影響を考慮した防護量も別の観測量である。
本章には先行の光量子比較1問と今回の現代物理8問を収める。原典と全候補を確認した解説初稿であり、平成23年度のβ安定への全称は科学的疑義を保持する。

\section{波と粒子のどちらを観測しているか}\label{theory:modern:duality}
量子論では光のエネルギー・運動量は \(E=hf\)、\(p=h/\lambda\) である。
光電効果では入射光子のエネルギーから仕事関数 \(W\) を差し引き、
\(K_{\max}=hf-W\) として光電子の最大運動エネルギーを説明する\cite{os_photoelectric}。
強度を上げることと一光子のエネルギーを上げることを分ける。
コンプトン散乱は光子と電子のエネルギー・運動量保存から波長変化を説明する例である。
一方、電子線回折は電子が運動量 \(p\) に対する物質波の波長 \(\lambda=h/p\) を持つことを示す。
結晶の周期構造による干渉は電子の波動性の観測であり、光の粒子性そのものの直接観測ではない\cite{os_matter}。
粒子名だけで「粒子性」と判断せず、現象が示す対象と性質を識別する。
平成24年度の比較問題をこの節へ移している。光電効果で変わる一光子エネルギーと、電子線の運動量に対応する干渉波長を混同しない。

\section{物質波：同速・同エネルギーは別の条件}\label{theory:modern:matter}
粒子の物質波長は \(\lambda=h/p\) である。同速vの自由粒子を比較すると \(p=\gamma mv\) の \(\gamma=(1-v^2/c^2)^{-1/2}\) が共通なので \(\lambda\propto1/m\) となる。電子、陽子、α核の静止質量はこの順に大きく、波長はこの順に長い。
非相対論的な同運動エネルギーKなら \(p=\sqrt{2mK}\) なので \(\lambda\propto m^{-1/2}\) となる。同速と同Kで質量の依存指数は異なる。同じ運動量で比較するなら波長は等しい。
電子線の回折から運動量と結晶構造を読む用途へつながる。電子の電荷の大きさだけでは波長順を求められない\cite{os_matter}。

\section{連続X線の最短端：一光子へのエネルギー上限}\label{theory:modern:xray}
初期運動エネルギーが無視できる電子を電圧Vで加速すれば、電場の仕事は \(eV\) である。金属標的での減速に伴う制動放射は連続スペクトルを持つが、一光子のエネルギーは電子が持つエネルギーを超えない。
\[
 hf_{\max}=eV,\qquad
 \lambda_{\min}=\frac{hc}{eV}.
\]
この上限はエネルギー保存と光子 \(E=hc/\lambda\) による。電子速度を \(mv^2/2\) から求める非相対論近似は不要である。金属の特性X線の離散線と連続線の端を区別する。
平成27年度は指定の \(h=6.6\times10^{-34}\)、\(e=1.6\times10^{-19}\)、\(c=3.0\times10^8\)、35kVを使う。現在の定義定数に勝手に置き換えず、約 \(3.54\times10^{-11}\,\mathrm m\) を得る。波長の長い光子が存在し得ることは、それを最短端の正答にする理由にならない\cite{os_photoelectric}。

\section{光素子：吸収・再結合・誘導放出}\label{theory:modern:devices}
半導体や分子の電子にはエネルギー状態がある。光子を吸収して高い状態へ移る過程と、励起された状態から低い状態へ移って光子を出す過程を使い分ける。
太陽電池は吸収した光によるキャリア生成と分離を使い、電流を取り出す。LEDは注入した電子・正孔の放射再結合を使う。有機ELは有機半導体の電子励起・励起子等の放射緩和を使う。全ての生成キャリアが光へ戻るわけではなく非放射過程も競合する\cite{modern_pv,modern_led,modern_oled}。
半導体レーザーでは、励起された活性層で誘導放出を増幅し共振器の帰還と組み合わせる。LEDの自発放出と、光場に相関した増幅を行うレーザーの機構を区別する。励起と光が関係するという同じ枠組みでも、素子の役割は同じではない\cite{modern_laser}。

\section{超伝導を光素子との相対比較で読む}\label{theory:modern:superconduct}
超伝導は臨界条件下で直流の無抵抗と特徴的な磁場応答を示す巨視的量子状態である。通常のBCS理論では電子対の凝縮を説明し、全ての超伝導体の微視的媒介機構を同じフォノン機構へ固定するわけではない。
電子が光を吸収・放出することを中心原理とする前節の4素子に対し、超伝導という現象の中心は凝縮状態の電気磁気応答である。令和5年度の比較はこの違いで判断する。超伝導体が電磁場と相互作用しない、光で状態を変えられないという断言へ広げない\cite{modern_superconduct,modern_pairs}。

\section{核種と放射線：A・Zの帳簿}\label{theory:modern:nuclear}
核種は陽子数Zと中性子数Nで指定し、A=Z+Nである。同位体は同Z・異Aであり、同質量数核とは異なる。
α放出では \((A,Z)\mapsto(A-4,Z-2)\)、βマイナスでは \((A,Z)\mapsto(A,Z+1)\)、βプラス又は電子捕獲では \((A,Z)\mapsto(A,Z-1)\) となる。γ脱励起はA・Zを保ち励起エネルギーを放射へ渡す。ニュートリノを含むエネルギー・運動量等の保存も必要である。
U・Thの主要天然壊変系列はα・βの連なりを経て安定Pb同位体へ至る。系列中でZが変わることと、個々のγ放出がZを変えないことは矛盾しない\cite{modern_mit_nuclear,modern_series}。

\section{β安定：質量最小という省略の限界}\label{theory:modern:stability}
同Aで質量をZの関数にすると、液滴模型のクーロン項と非対称項からおおよそ放物線を得る。ただし核の偶奇に伴う対エネルギーがあり、偶数Aでは偶偶核と奇奇核の曲線が異なる。
β壊変は隣接Zの終状態へQ値が正となる場合にエネルギー的に可能になる。原子質量を用いたβマイナスは、電子束縛等を省く近似で \(Q_{\beta^-}=[M(A,Z)-M(A,Z+1)]c^2\) である。βプラスの原子質量式には \(2m_ec^2\) のしきい値がある。可能性と寿命は核状態の選択則等も含み同じではない。
\begin{tcolorbox}[title=平成23年度の最小質量への全称を保持しない]
偶数Aでは単一βに対して安定な核が複数あり得る。中間の奇奇核が高い質量にあるため、一つのβ遷移で最小質量へ進めない場合がある。二重β壊変の観測はあるが、その経路・極めて長い時間尺度は単一β連鎖と同一ではない。「全ての同A核が最小質量核へβ壊変する」という原典の省略を、無条件の定理として教材へ移さない。
\end{tcolorbox}
公式2という熱中性子の取り違えは独立に説明できるが、3の一般傾向と全称の科学的一意性を分けて査読へ送る\cite{modern_mit_mass,modern_doublebeta,modern_mit_nuclear}。

\section{核分裂性と中性子のエネルギー}\label{theory:modern:fission}
重い核が中性子を吸収し、複合核の励起が分裂障壁等に対して有効となると核分裂が起こり得る。核種によって必要な中性子エネルギーが違う。
235U等は熱中性子で核分裂しやすい核分裂性核種である。238Uは高速中性子による核分裂が可能で、熱中性子の捕獲からβ壊変を経てPuへつながる親物質でもある。
熱中性子炉の連鎖反応で、自然存在量の多い238Uをそのまま熱核分裂の主役として説明しない。核分裂可能と熱核分裂性、単発反応と持続連鎖を分ける\cite{modern_fastfission}。

\section{壊変率と半減期：何が減った割合か}\label{theory:modern:decay}
同じ核種の単位時間の壊変確率を一定の \(\lambda\) とし、生成・流出入がなければ
\[
 \frac{dN}{dt}=-\lambda N,\qquad
 N(t)=N_0e^{-\lambda t},\qquad
 T_{1/2}=\frac{\ln2}{\lambda},\qquad
 \frac{N(t)}{N_0}=2^{-t/T_{1/2}}.
\]
半減期は一つの核が必ずその時刻に壊変する時間ではなく、集団の平均残存が半分になる時間である。放射能 \(\mathcal A=\lambda N\) も同じ核種なら同じ比で減る。
90年/30年は3半減期で1/8、30日/8日は約0.0743、30日/9日は約0.0992である。整数の半減期数に丸めず指数で評価する。元の親核種の残存比と壊変済み割合 \(1-N/N_0\)、娘核種量を区別する。娘も壊変する系列では、娘量はこの単純な差だけではない\cite{modern_mit_nuclear}。

\section{吸収線量・等価線量・実効線量}\label{theory:modern:dose}\label{theory:modern:activity}
吸収線量Dは単位質量当たりに吸収された放射線エネルギーで、Gy=J/kgである。Bqは壊変率であり、核種の活動度だけからそのままGyへ換算できない。
組織Tに放射線Rが与える平均吸収線量 \(D_{T,R}\) に放射線の種類等の無次元重み \(w_R\) を付け、
\[
 H_T=\sum_Rw_RD_{T,R},\qquad E=\sum_Tw_TH_T
\]
を等価線量と実効線量とする。どちらもSvである。実効線量の \(w_T\) は組織の寄与を重み付ける。原典の「荷重係数」という表記を校正で別の語へ置換せず、重みとして意味を説明する。
GyとSvは同じJ/kgの次元であるが量の定義が違う。等価線量・実効線量は防護上の量であり、高線量の確定的影響や個人への治療判断をそれだけで表す量へ転用しない。
α線は外部透過が弱くても、内部取込みでは高い密度でエネルギーを与え得る。実際の影響には量・体内分布・経路等が必要である。「紙で止まるから常に影響なし」「内部なら常に同じ被害」のどちらも単位の定義からは導けない\cite{modern_dose,modern_alpha}。

\section{出題形式と理解の接続}
今回の8問は、同速の物質波1、制動X線端1、半減期3、核反応比較1、線量対応1、光電子素子比較1である。先行の光量子比較を加え9問を収録する。これらは確認済みの当該問題集合の集計であり将来出題の保証ではない。
エネルギー上限、運動量と波長、保存するA・Z、指数減衰、測定量の意味の順に整理すると、式を暗記するだけでなく条件の変更へ転用できる。地球の同位体年代では親娘閉鎖と試料反応がさらに必要であり、本章の壊変式だけで地質年代を確定しない。

\section{過去問と全候補の検討}
\subsection{平成23年度 問10}
\label{q:h23-reg-applied-10}

\noindent\textbf{原典：} sources/H23/17\_applied\_science.pdf、PDF物理4頁。

\textbf{公式正答：} 2（single）。

\par\textit{原典本文の校正と独自解説は未作成。分類のみの登録項目。}

\subsection{平成25年度 問7}
\label{q:h25-reg-applied-07}

\noindent\textbf{原典：} sources/H25/17\_applied\_science.pdf、PDF物理4頁。

\textbf{公式正答：} 5（single）。

\par\textit{原典本文の校正と独自解説は未作成。分類のみの登録項目。}

\subsection{平成26年度 問9}
\label{q:h26-reg-applied-09}

\noindent\textbf{原典：} sources/H26/17\_applied\_science.pdf、PDF物理4頁。

\textbf{公式正答：} 2（single）。

\par\textit{原典本文の校正と独自解説は未作成。分類のみの登録項目。}

\subsection{平成27年度 問5}
\label{q:h27-reg-applied-05}

\noindent\textbf{原典：} sources/H27/17\_applied\_science.pdf、PDF物理2頁。

\textbf{公式正答：} 2（single）。

\par\textit{原典本文の校正と独自解説は未作成。分類のみの登録項目。}

\subsection{平成27年度 問6}
\label{q:h27-reg-applied-06}

\noindent\textbf{原典：} sources/H27/17\_applied\_science.pdf、PDF物理3頁。

\textbf{公式正答：} 5（single）。

\par\textit{原典本文の校正と独自解説は未作成。分類のみの登録項目。}

\subsection{平成28年度 問9}
\label{q:h28-reg-applied-09}

\noindent\textbf{原典：} sources/H28/17\_applied\_science.pdf、PDF物理4頁。

\textbf{公式正答：} 3（single）。

\par\textit{原典本文の校正と独自解説は未作成。分類のみの登録項目。}

\subsection{令和4年度 問9}
\label{q:r04-reg-applied-09}

\noindent\textbf{原典：} sources/R04/17\_applied\_science.pdf、PDF物理3頁。

\textbf{公式正答：} 2（single）。

\par\textit{原典本文の校正と独自解説は未作成。分類のみの登録項目。}

\subsection{令和5年度 問7}
\label{q:r05-reg-applied-07}

\noindent\textbf{原典：} sources/R05/17\_applied\_science.pdf、PDF物理3頁。

\textbf{公式正答：} 4（single）。

\par\textit{原典本文の校正と独自解説は未作成。分類のみの登録項目。}

